\originalchapter{空间的拼接与基本群}\label{ch:vankampen}
\emph{van Kampen定理及曲面基本群与18.904的教学要求对应; 胞腔粘合的模型参照18.905第14讲. 本章自由积、网格论证、附加关系和完整题解为编者补充.}

\originalsection{自由群与自由积}
一个环路在不同开集之间穿行时, 可以分段记录每段的基本群元素. 分段记录首先形成一个词, 再由交集中的相容关系把不同记录识别.
\begin{definition}
集合$S$上的自由群$F(S)$由符号$s,s^{-1}$组成的有限约化词构成, 约化指没有相邻的$s s^{-1}$或$s^{-1}s$. 空词为单位, 乘法为拼接后消去相邻逆符号. 两群$A,B$的自由积$A*B$由来自两群的非单位元素交替组成的词构成, 同一因子的相邻元素合并为其乘积, 乘积为单位时删去.
\end{definition}
约化结果唯一可用逐字扫描说明. 用一串已经约化的符号作栈, 新符号若与栈顶相消就删去栈顶, 否则加入. 对自由积, 同因子栈顶与新元素先在因子中相乘, 单位即弹出, 非单位即替换. 在任意位置先合并或删去一个允许的相邻对, 再扫描, 与原词扫描结果一致; 这一事实对自由积只用因子内部的结合律, 对自由群只用逆符号消去. 因而任何消去次序都得相同的约化词. 特别$R(R(u)v)=R(uv)=R(uR(v))$, 其中$R$表示约化, 故乘法结合. 反向词并逐字取逆给逆元.

\begin{proposition}\label{thm:free-universal}
任意$S\to G$唯一扩充为$F(S)\to G$的同态. 给定同态$A\to G$, $B\to G$, 唯一得到$A*B\to G$的同态.
\end{proposition}
\begin{proof}
把词的每个符号映到指定元素, 按次序相乘. 相邻逆元的消去或同因子乘积的合并不改变结果, 故在约化词上良定义. 拼接对应乘法, 故为同态. 每个词由生成元素相乘构成, 唯一性随之成立.
\end{proof}
给群$G$的一组元素$R$, 记$\langle\!\langle R\rangle\!\rangle$为包含它们的最小正规子群, 即所有共轭$g r^{\pm1}g^{-1}$的有限乘积. 商群$G/\langle\!\langle R\rangle\!\rangle$将这些元素强制等于单位. 记号$\langle S\mid R\rangle$表示从$F(S)$添加关系$R$所得的群. 例如$\langle a,b\mid aba^{-1}b^{-1}\rangle\cong\mathbb Z^2$: 关系允许$a,b$交换, 每个词化成$a^m b^n$, 映到$(m,n)$表明不会产生别的相等关系.

\originalsection{van Kampen定理}
\begin{theorem}\label{thm:vankampen}
设$X=U\cup V$, $U,V$开且道路连通, $W=U\cap V$非空且道路连通, $x_0\in W$. 记包含诱导映射$i_U:\pi_1(W,x_0)\to\pi_1(U,x_0)$, $i_V:\pi_1(W,x_0)\to\pi_1(V,x_0)$. 则自然同态给出
\[
 \pi_1(X,x_0)\cong
 \frac{\pi_1(U,x_0)*\pi_1(V,x_0)}
 {\langle\!\langle i_U(w)i_V(w)^{-1}:w\in\pi_1(W,x_0)\rangle\!\rangle}.
\]
不要求$i_U,i_V$单射.
\end{theorem}
\begin{proof}
记右侧为$G$, 两因子到$G$的同态为$j_U,j_V$. 它们对来自$W$的环路给出相同元素. 包含到$X$的同态也有这种相容性, 因而自由积泛性质给出$\Phi:G\to\pi_1(X,x_0)$.

要构造逆映射, 对每个$z\in X$固定一条$c_z$从$x_0$到$z$: 当$z\in W$时选在$W$中, 当$z$只在$U$或只在$V$时选在相应开集中, 并令$c_{x_0}$常值. 若一段道路$\gamma$的像全在$U$, 给它赋值$j_U[c_{\gamma(0)}*\gamma*\overline{c_{\gamma(1)}}]$; 全在$V$时用$j_V$. 若同时在二者, 整个环路在$W$内, 两种赋值相同. 反向道路赋逆元素.

对环路$\alpha$, 将$I$细分使每段像落在$U$或$V$, 定义$\Theta(\alpha)$为这些赋值按次序的乘积. 在同一区域内加一个分点, 中间出现的$\overline c*c$可消去, 所以乘积不变. 两组分割有共同细分; 同一小段若被不同区域记录, 像必在$W$, 赋值仍相同. 因而$\Theta$与分割无关.

还须证明它与同伦无关, 这正是关系完备性的关键. 对固定基点同伦$H:I^2\to X$, 用$H^{-1}(U),H^{-1}(V)$的Lebesgue数将正方形分成有限小格, 每格像落在一个区域. 赋值于各有向格边. 单个格子的边界在其区域内可缩, 因为先将格边界在凸格内缩到一角, 再与$H$复合; 加上的$c_z$与逆道路互相消去. 因而每格都满足“下边乘右边等于左边乘上边”的群等式. 逐行逐格将一条网格道路移过小格, 每次使用这个等式, 得整个正方形下边乘右边等于左边乘上边. 两侧边在$x_0$恒定, 赋值为单位. 所以底环路与顶环路赋值相同.

由此$\Theta:\pi_1(X,x_0)\to G$良定义, 拼接分割表明它为同态. 对全在$U$或$V$内的环路, 它就是$j_U$或$j_V$. 因而$\Theta\Phi$在两个生成因子上为恒等, 在$G$上为恒等. 另一方面,$\Phi\Theta[\alpha]$把各段连同$c_z$拼接回$X$, 相邻逆道路消去后就是$[\alpha]$. 所以两映射互逆.
\end{proof}
开集条件保证有限细分和网格均存在, 交集道路连通保证可以在重叠处选择共同连接道路. 若交集不连通, 该版本不能直接应用, 即使图上两个区域看起来都很简单.

\originalsection{图与球面}
\begin{proposition}\label{thm:graph-pi1}
$k$个圆周粘在一点的楔和的基本群是$k$个生成元的自由群. 有限连通图有$V$个顶点、$E$条边时, 基本群是秩$E-V+1$的自由群.
\end{proposition}
\begin{proof}
对楔和, 取$U$为前$k-1$个圆周及第$k$圈在粘合点附近的短开弧, $V$为第$k$圈及其余圈在粘合点附近的短弧. 二者相对开, 分别形变收缩到前$k-1$圈与第$k$圈, 交集是一棵星状树, 可缩. van Kampen给出自由积, 由$\pi_1(S^1)=\mathbb Z$归纳.

一般连通图先选生成树$T$. 它有$V-1$条边: 从一个顶点起, 每次加入通向新顶点的边, 不造回路, 最后遍历全部顶点. 树可缩, 由不断将叶边线性压向其另一端得到. 再逐条加入其余边. 加入边的两端在$T$中由唯一简单道路连接. 对新图取$U$为去掉新边中点的空间, 它沿新边两侧开尾形变收缩到旧图; 取$V$为新边与该树中道路的一个小开邻域, 它形变收缩到所构成的圆周. 邻域在道路以外只含各岔边的短开尾, 这些尾均可线性压到所附顶点. $U\cap V$为删掉新边中点后的上述圆周连同短开尾, 可缩且道路连通. 故每加一条非树边, 基本群自由积上一个$\mathbb Z$. 一共$E-(V-1)$条, 得所述秩. 环边和重复边按独立边处理, 论证仍成立.
\end{proof}

\begin{proposition}\label{thm:sphere-simply-connected}
对$n\ge2$, $S^n$单连通. 对$n\ge2$, $\pi_1(\mathbb RP^n)=\mathbb Z/2\mathbb Z$.
\end{proposition}
\begin{proof}
以去掉南极、去掉北极的两个开集覆盖$S^n$, 各自经立体投影同胚于$\R^n$, 故单连通. 交集同胚于$S^{n-1}\times\R$, 在$n\ge2$时道路连通. van Kampen右侧为平凡群的自由积再取商, 仍平凡, 所以$S^n$单连通.

对径商$p:S^n\to\mathbb RP^n$为两层覆盖. 对点$x$取足够小开球帽$U$使$U\cap(-U)=\varnothing$, 饱和开集$U\cup(-U)$投影为开集, 且两个球帽各同胚到它; 此处用商判据及对径作用的开性. $S^n$单连通, 因而这是普适覆盖, 两层使底空间基本群只有两个元素, 即$\mathbb Z/2\mathbb Z$. 或用甲板群: 只有恒等与对径变换.
\end{proof}
当$n=1$时$\mathbb RP^1$本身为圆周, 基本群为$\mathbb Z$, 不能删掉$n\ge2$条件.

\originalsection{附加一个二维胞腔}
\begin{theorem}\label{thm:attach-cell-pi1}
设$A$为道路连通的紧Hausdorff空间, 沿连续$f:S^1\to A$附加一个圆盘, 得$X=A\cup_f D^2$. 以$f(1)$为基点, 则
$\pi_1(X)\cong\pi_1(A)/\langle\!\langle[f]\rangle\!\rangle$.
有限个圆盘依次附加时, 对全部附加环路取正规闭包.
\end{theorem}
\begin{proof}
在圆盘上以$r=|z|$为径向参数. 取$U$为$A$连同圆盘的$r>1/3$部分, $V$为圆盘内部$r<2/3$部分. 它们在商中开: 在不交并$A\sqcup D^2$中的饱和逆像开, 可用商判据. $U$将环带径向推到边界后形变收缩到$A$, $V$为开圆盘而可缩, 交集为开环带. 对$U$的形变在$A$上恒等、圆盘上沿半径移动, 并在边界与$f$一致. 为严格验证时间参数中的连续性, 注意$A\sqcup D^2$紧Hausdorff, 附加的等价关系闭: 它由对角线、边界与$f$的图像及其转置、满足$f(z)=f(z')$的边界点对组成, 各部分均闭. 定理\ref{thm:closed-equivalence}使$X$紧Hausdorff. 因而商映射与$I$之积也是从紧空间到Hausdorff空间的连续满射, 为商映射; 限制到开集$U\times I$仍为商映射. 径向同伦在不交并的相应开子集与$I$之积上连续并在纤维中常值, 可由商泛性质下降到$U\times I$.

交集的基本群由绕中心一圈生成, 包含到$U$并收缩到$A$后正是附加环路$f$, 包含到$V$则为单位. 用van Kampen得到所述正规闭包商. 基点若先取在交集中, 用径向道路改变到$f(1)$, 只改变附加环路的共轭代表, 正规闭包不变. 多个圆盘逐次重复此步骤.
\end{proof}
这个定理用于有限图和有限胞腔的附加, 下面的具体空间均满足所用的商邻域条件. 对更一般粘合, 不能将任意闭子空间的并直接当作本定理中的开覆盖.

\begin{example}
计算由边界词$aba^{-1}b^{-1}$粘合正方形所得环面, 以及由$a^2$附加圆盘所得实射影平面的基本群.
\end{example}
\begin{solution}
环面的一维骨架是两个圆周的楔和, 基本群$F(a,b)$. 填入面使边界词成为单位, 得$\langle a,b\mid aba^{-1}b^{-1}\rangle=\mathbb Z^2$, 与积空间计算一致. 射影平面的圆盘边界将对径点识别, 边界映到射影直线时绕行两次, 故附加词为$a^2$. 得$\langle a\mid a^2\rangle=\mathbb Z/2\mathbb Z$, 与覆盖计算一致.
\end{solution}

\originalsection{练习与解答}
\begin{exercise}
证明$\R^2\setminus\{(-1,0),(1,0)\}$的基本群是两个生成元的自由群. 给出绕两个点的绕数都为零但不可缩的环路类.
\end{exercise}
\begin{solution}
取$U=\{x<1/2\}\setminus\{(-1,0)\}$, $V=\{x>-1/2\}\setminus\{(1,0)\}$. 每个穿孔半平面同胚于穿孔平面, 基本群为$\mathbb Z$; 交集为竖直开条带而凸. van Kampen给出$F(a,b)$, $a,b$分别绕两点. 交换子$aba^{-1}b^{-1}$约化后仍是非空词, 所以不可缩. 它对每个点的总绕数均为零, 因为对应同态分别只记录$a,b$的指数和. 因而两个普通绕数不能检测非交换信息.
\end{solution}
\begin{exercise}
空间由一个圆周及沿$z\mapsto z^m$, $z\mapsto z^n$附加的两个圆盘组成, 其中$m,n\ge1$. 求基本群.
\end{exercise}
\begin{solution}
得到$\langle a\mid a^m,a^n\rangle$. 原圆周群为$\mathbb Z$, 被添加关系生成的子群为$m\mathbb Z+n\mathbb Z=d\mathbb Z$, $d=\gcd(m,n)$. 一个方向所有组合均为$d$倍数; 反向B\'ezout等式$um+vn=d$使$d$也在生成子群中. 故基本群为$\mathbb Z/d\mathbb Z$; 特别互素时单连通, 但不能据此断言空间可缩.
\end{solution}
\begin{exercise}
证明Klein瓶的边界词$aba^{-1}b$给出群$\langle a,b\mid aba^{-1}=b^{-1}\rangle$, 并证明此群不是交换群.
\end{exercise}
\begin{solution}
由附加胞腔定理得到该表示. 在$\R^2$上定义$A(x,y)=(x+1,-y)$, $B(x,y)=(x,y+1)$, 则$ABA^{-1}=B^{-1}$. 因而泛性质给出此表示群到这些变换群的同态. 两变换$AB$和$BA$分别将$(0,0)$送到$(1,-1)$和$(1,1)$, 不同, 所以群中$a,b$也不交换. 这一步用具体作用检验非交换性, 而不是仅看关系的外形.
\end{solution}
